% =============================================================
%  第七单元  乘法口诀（二）
% =============================================================
\Chapter{第七关\quad 八九七十二的来历}
\noindent{\small\kaishu\color{tipc}（课本第七单元：乘法口诀（二），P77--P86）}

\begin{blunder}
豆角爸爸三十多岁，还要掰手指算 $9\times7$。被小豆角抓了个正着。
小豆角说：“豆角爸爸不是大人吗？”
从那以后小豆角每天考豆角爸爸一句，豆角爸爸只好偷偷练了两天。
\textbf{所以这一关豆角爸爸跟小豆角一起背}——不是监督小豆角，是豆角爸爸也得会。
\end{blunder}

\begin{gamble}{赌一把：9 的口诀手指赛}
\begin{itemize}
  \item 两人同时开始，从 $9\times1$ 报到 $9\times9$，用弯手指的办法，谁快谁赢。
  \item 第二局换个玩法：不许用手指，只许用“$10n-n$”算。\textbf{这局小豆角多半输}——
        正好让小豆角明白：手指法快，但换个办法更厉害。
  \item 输的人\textbf{给对方一块钱零花钱}。
\end{itemize}
\end{gamble}

\begin{bonus}
\begin{itemize}
  \item \textbf{西游记接龙}：豆角爸爸起一句“八九七十二变”，小豆角接“七七四十九天”，
        再接“九九八十一难”。接不上的人负责收碗。
  \item \textbf{一天一句}：每天上学路上随口抽一句（跳着抽，别按顺序），
        答不上来也不批评，第二天再抽同一句。
  \item \textbf{考豆角爸爸}：让小豆角随便挑五句考豆角爸爸，豆角爸爸\textbf{故意错两道}。
        小豆角每天都会来考豆角爸爸——这是最好的复习。
\end{itemize}
\end{bonus}

\begin{unitinfo}
\textbf{课本 P77--P86}\quad \textbf{6 次课，每次 20 分钟}\quad
\textbf{建议第 18--20 周} \\[4pt]
\textbf{这一关到底在干嘛}：把第四单元那个“自己编口诀”的动作，
再来四遍（6、7、8、9），然后收成一张 45 句的表。\\
\textbf{逻辑是这么走的}：\ 编口诀（照老办法）$\to$ 忘了一句就\textbf{拆开算}
$\to$ 找规律（相邻差几）$\to$ 收成大表 $\to$ 横着会、竖着会、跳着会。\\
\textbf{这一章背诵量最大，但也最不该硬背。}
判断小豆角该不该回头补：问小豆角“三七二十一这句是怎么来的？”
答不出来，说明第四单元是死记的，\textbf{必须回去补编一次}，不然这四组会全乱。
\end{unitinfo}

\subsection{开局先聊个事：这些数生活里到处都是}

一盒鸡蛋，一排 6 个；一周 7 天；九宫格 9 格；
一层楼梯 9 级；一箱牛奶一排 8 个。

到了这一章，“几个几”已经不用再解释了，小豆角手上有数。
所以这一章的重点不再是“懂”，而是\textbf{把 45 句变成能随时调出来的工具}。

\subsection{任务 1：编 6 的口诀}

\begin{lesson}{1}{20}
\textbf{今天要拿下的}：照第四单元的办法，自己编出 6 的口诀。

\begin{materials}
百数表、笔、纸。
\end{materials}

\begin{enumerate}
  \item \textbf{给情境}（课本 P77）：每行贴 6 张贴画。
        \begin{center}
        \small
        \begin{tabular}{@{}l*{9}{c}@{}}
        \toprule
        有几行 & 1 & 2 & 3 & 4 & 5 & 6 & 7 & 8 & 9 \\
        \midrule
        有几张 & 6 & 12 & 18 & 24 & 30 & 36 & 42 & 48 & 54 \\
        \bottomrule
        \end{tabular}
        \end{center}
        让小豆角自己填（每次加 6）。
  \item \textbf{让小豆角自己编}（豆角爸爸不出声）：一六得六、二六十二、三六十八、
        四六二十四、五六三十、六六三十六、六七四十二、六八四十八、六九五十四。
  \item \textbf{找规律}：相邻两句差几？（6）得数有什么共同点？（全是双数）
  \item \textbf{找锚点}：只要牢牢记住 \textbf{六六三十六}，
        往后就是 36、42、48、54——每次加 6，不用另记。
        \textbf{这就是“少背”的窍门：先记一个点，再靠加。}
  \item \textbf{回头对一遍}：6 的口诀里，哪句和前面的重复？
        “二六十二”和“三四十二”都得 12。让小豆角自己找，找不到就留到学完 8 再回头找。
\end{enumerate}

\begin{passbox}
\begin{itemize}
  \item 6 的口诀能独立编出来。
  \item 小豆角能说出“相邻两句差 6”，并知道记住“六六三十六”就能往后推。
\end{itemize}
\end{passbox}
\end{lesson}

\subsection{任务 2：忘了就拆开算}

\begin{lesson}{2}{20}
\textbf{今天要拿下的}：学会把忘掉的那句\textbf{拆成两句会的加起来}。

\begin{materials}
方格纸（画点子图）、彩笔。
\end{materials}

\begin{enumerate}
  \item \textbf{先问}：“$6\times7$ 小豆角要是忘了，怎么办？”（别说“再背一遍”）
  \item \textbf{拆法一：拆成 5 和 2}：
        $6\times7 = 6\times5 + 6\times2 = 30+12 = 42$。
        \begin{center}
        \begin{tikzpicture}[scale=0.42]
          \foreach \x in {0,...,6}{\foreach \y in {0,...,5}{
            \fill[mainc!25] (\x,\y) circle (2.2pt);}}
          \foreach \x in {0,...,4}{\foreach \y in {0,...,5}{
            \fill[mainc] (\x,\y) circle (2.2pt);}}
          \draw[pitc,thick] (-0.32,-0.32) rectangle (4.32,5.32);
          \node[pitc,below,font=\scriptsize] at (2,-0.45) {$6\times5=30$};
          \node[sayc,below,font=\scriptsize] at (5.5,-0.45) {$6\times2=12$};
        \end{tikzpicture}
        \end{center}
        （6 行 7 列，先把左边 5 列框出来。）
  \item \textbf{拆法二：拆成 6 和 1}：$6\times7 = 6\times6 + 6\times1 = 36+6=42$。
  \item \textbf{拆法三：反过来}：$6\times7 = 7\times6$（等 7 的口诀学完就能用）。
  \item \textbf{做四道拆数题，先画图再算}：
        \begin{enumerate}[label=（\arabic*）,itemsep=2pt]
          \item $6\times8 = 6\times5 + 6\times$ \blank[1cm] $=$ \blank[1.4cm]
          \item $6\times9 = 6\times$ \blank[1cm] $+ 6\times$ \blank[1cm] $=$ \blank[1.4cm]
          \item $8\times8 = 8\times$ \blank[1cm] $+ 8\times$ \blank[1cm] $=$ \blank[1.4cm]
          \item $7\times7 = 7\times$ \blank[1cm] $- 7\times$ \blank[1cm] $=$ \blank[1.4cm]
        \end{enumerate}
        答案：（1）3；48\quad（2）如 $6\times5+6\times4=54$\quad
        （3）如 $8\times4+8\times4=64$\quad（4）如 $7\times8-7\times1=49$。
  \item \textbf{收一句贴桌上}：“\textbf{忘了就拆，拆成两句会的加起来。}”
        \textbf{会拆的小豆角永远不会被一句口诀卡死。}
\end{enumerate}

\begin{saybox}
“$6\times7$ 忘了，小豆角能用会的算出来吗？怎么拆？”\\
“为什么可以拆？”（7 列可以分成 5 列加 2 列，点子一个没少）\\
“第 4 题为什么要减？”（多算了 1 个 7，得减掉）
\end{saybox}

\begin{passbox}
\begin{itemize}
  \item 能用两种不同的拆法算出 $6\times7$。
  \item 四道拆数题全对，都能画出图说明。
\end{itemize}
\end{passbox}
\end{lesson}

\subsection{任务 3：7、8、9 的口诀（西游记里全是口诀）}

\begin{lesson}{3}{20}
\textbf{今天要拿下的}：编出 7、8、9；用《西游记》的故事帮着记。

\begin{materials}
百数表、笔、纸。
\end{materials}

\begin{enumerate}
  \item \textbf{先讲三个故事}（课本 P79）：
        \begin{itemize}
          \item 孙悟空有\textbf{八九七十二}般变化；
          \item 在炼丹炉里炼了\textbf{七七四十九}天，炼出火眼金睛；
          \item 唐僧取经经历\textbf{九九八十一}难。
        \end{itemize}
        \textbf{先给这三句}，让小豆角觉得口诀是有来处的，不是凭空一堆字。
  \item \textbf{编 7}（每次加 7）：一七得七、二七十四、三七二十一、四七二十八、
        五七三十五、六七四十二、七七四十九、七八五十六、七九六十三。
  \item \textbf{编 8}（每次加 8）：一八得八、二八十六、三八二十四、四八三十二、
        五八四十、六八四十八、七八五十六、八八六十四、八九七十二。
  \item \textbf{编 9}（每次加 9）：一九得九、二九十八、三九二十七、四九三十六、
        五九四十五、六九五十四、七九六十三、八九七十二、九九八十一。
  \item \textbf{交叉找“双胞胎”}（这一步能省掉一半力气）：
        \begin{center}
        \small
        \begin{tabular}{@{}ll@{}}
        \toprule
        \textbf{得数} & \textbf{对应口诀} \\
        \midrule
        12 & 二六十二、三四十二（两句） \\
        18 & 三六十八、二九十八（两句） \\
        24 & 四六二十四、三八二十四（两句） \\
        36 & 六六三十六、四九三十六（两句） \\
        42 & 六七四十二 \quad 48：六八四十八 \\
        54 & 六九五十四 \quad 56：七八五十六 \\
        63 & 七九六十三 \quad 72：八九七十二 \\
        \bottomrule
        \end{tabular}
        \end{center}
        \textbf{重点盯那四组“双胞胎”}（12、18、24、36），考试最容易混的就是它们。
  \item \textbf{背的顺序别从 1 背到 9}：先背\textbf{平方数}
        （六六三十六、七七四十九、八八六十四、九九八十一），
        它们是锚点，其余靠加减一个推。
\end{enumerate}

\begin{saybox}
“孙悟空会几变？”（七十二变 $\to$ 八九七十二）\\
“哪个得数有两句口诀？找出三组。”\\
“小豆角觉得先记住哪几句最省事？”（平方数）
\end{saybox}

\begin{passbox}
\begin{itemize}
  \item 7、8、9 的口诀能独立编出来（可看表）。
  \item 能找出至少三组“得数相同的两句口诀”。
\end{itemize}
\end{passbox}
\end{lesson}

\subsection{任务 4：9 的口诀，用十根手指}

\begin{lesson}{4}{20}
\textbf{今天要拿下的}：用手指记住 9 的口诀；理解 $9\times n = 10n-n$。

\begin{materials}
小豆角自己的两只手。不用别的。
\end{materials}

\begin{enumerate}
  \item \textbf{手指法}（课本 P80）：
        \begin{center}
        \begin{tikzpicture}[x=0.62cm,y=1cm,font=\scriptsize]
          \foreach \i in {1,...,10}{
            \pgfmathsetmacro{\xx}{\i-1}
            \ifnum\i=4
              \draw[pitc,thick] (\xx,0) -- (\xx,0.45);
              \draw[pitc,thick] (\xx,0.95) -- (\xx,1.4);
            \else
              \draw[tipc,thick] (\xx,0) -- (\xx,1.4);
            \fi
            \node[below] at (\xx,-0.08) {\i};
          }
          \node[pitc,above=2mm] at (3,1.4) {弯下第 4 根};
          \node[mainc] at (1,1.75) {左边 3 根 $\to$ 30};
          \node[mainc] at (6.2,1.75) {右边 6 根 $\to$ 6};
          \node[tipc] at (4.6,2.35) {\textbf{$9\times4=36$}};
        \end{tikzpicture}
        \end{center}
        \textbf{做法}：伸出十指。算 $9\times n$，\textbf{从左往右数，弯下第 $n$ 根}。
        弯下去那根\textbf{左边剩几根，就是十位}；\textbf{右边剩几根，就是个位}。
  \item \textbf{当场验四次}：$9\times2$（弯第 2 根 $\to$ 1 和 8 $\to$ 18）、
        $9\times6$（$\to$ 5 和 4 $\to$ 54）、$9\times7$（$\to$ 63）、$9\times9$（$\to$ 81）。
        \textbf{让小豆角自己弯、自己数。}
  \item \textbf{更好用的规律}：$9\times n = 10n-n$。
        $9\times1=10-1=9$；$9\times2=20-2=18$；$9\times3=30-3=27$……
        让小豆角填完 $9\times5$ 到 $9\times9$：$45$、$54$、$63$、$72$、$81$。
  \item \textbf{第三个规律}：9 的口诀，\textbf{十位加个位都等于 9}
        （18、27、36、45、54、63、72、81）。而且十位 1 到 8 往上走，
        个位 8 到 1 往下走。\textbf{三个规律对三遍，9 的口诀基本不用背了。}
  \item \textbf{往表外推}（课本 P80 试一试）：$9\times10=?$ 用 $90-10=80$；
        $9\times11=9\times10+9=99$。
        \textbf{这说明口诀的规律是能一直往后用的}，不是死表。
\end{enumerate}

\begin{saybox}
“弯第 6 根，左边几根右边几根？”\\
“$9\times n$ 除了手指，还能怎么算？”（$10n-n$）\\
“9 的口诀得数，十位加个位等于几？”（全是 9）
\end{saybox}

\begin{pitbox}
\textbf{手指法的两个坑}：
\begin{enumerate}[label=\arabic*、]
  \item \textbf{从右手开始数}。必须\textbf{从左往右}（左手拇指是第 1 根）。
  \item \textbf{数错边}。提醒：“弯指\textbf{左边当十，右边当个}。”
        练的时候让小豆角\textbf{小声念出来}：“左边 3，30；右边 6，6。”
\end{enumerate}
\end{pitbox}

\begin{passbox}
\begin{itemize}
  \item 手指法能当场算出 9 的任意一句，左右不分错。
  \item 能说出“十位加个位都等于 9”。
  \item 能算出 $9\times10$、$9\times11$。
\end{itemize}
\end{passbox}
\end{lesson}

\subsection{任务 5：收成一张乘法表}

\begin{lesson}{5}{20}
\textbf{今天要拿下的}：45 句收成一张表，并且能从表里看出规律。

\begin{materials}
附录 B 的乘法表（填空版）、彩笔。
\end{materials}

\begin{enumerate}
  \item \textbf{填表}：用填空版，把乘数和积都填上。
        \textbf{这一课主要工作就在这儿，慢点没关系。}
  \item \textbf{一条一条看规律}（课本 P84 的问题）：
        \begin{itemize}
          \item \textbf{对角线}上是什么？（$1\times1$、$2\times2$……$9\times9$，全是平方数）
          \item \textbf{对称吗}？$3\times7$ 和 $7\times3$ 在哪？得数一样吗？（关于对角线对称）
          \item 每一\textbf{行}有什么规律？（第 3 行每次加 3，第 7 行每次加 7）
          \item 每一\textbf{列}呢？
          \item 哪几个得数出现两次？（12、18、24、36）
        \end{itemize}
  \item \textbf{找积是 18 的算式}（课本 P83）：$2\times9$、$3\times6$、$6\times3$、$9\times2$，
        四个算式、两句口诀。
  \item \textbf{找乘数里有 3 的算式}：$3\times1$ 到 $3\times9$，加 $1\times3$ 到 $9\times3$。
  \item \textbf{用表查}：随便报一个数（比如 56），让小豆角在表里找到位置，说出哪句口诀。
        \textbf{这一步是把表当工具用，不是当图片看。}
\end{enumerate}

\begin{saybox}
“表里哪一句最难记？单独做张卡片。”\\
“$3\times7$ 和 $7\times3$ 关于什么对称？”\\
“56 是哪句口诀？在表里指给豆角爸爸看。”
\end{saybox}

\begin{passbox}
\begin{itemize}
  \item 填空版乘法表独立填完，对 90\% 以上。
  \item 能说出三条规律（对角线、对称、每行加几）。
\end{itemize}
\end{passbox}
\end{lesson}

\subsection{任务 6：口诀总过关}

\begin{lesson}{6}{20}
\textbf{今天要拿下的}：横着、竖着、跳着都能答；给一个积能反过来说出乘数。

\begin{materials}
附录 A 的乘法卡、错题卡片。
\end{materials}

\begin{enumerate}
  \item \textbf{三遍法}（每遍 3 分钟，中间歇 1 分钟）：
        \begin{enumerate}[label=第\arabic*遍：]
          \item \textbf{横着}：1 的口诀到 9 的口诀，一行一行过。
          \item \textbf{竖着}：一一得一、一二得二、一三得三……到一九得九；
                再二二得四、二三得六……（按列走，难度高一档）
          \item \textbf{跳着}：豆角爸爸随便报（七八、四九、三六、五八），小豆角报得数。
        \end{enumerate}
  \item \textbf{错的做卡}：跳着背时答错的 3--5 句写卡片，贴镜子或书桌。
        \textbf{只贴错的，不贴整表。}
  \item \textbf{60 秒速算}：抽卡，60 秒对 30 张算过。
  \item \textbf{反着用}：报一个积（42）问“42 除以 6 是几？”$\to$ 想“六七四十二” $\to$ 7。
        \textbf{这是第八单元的前置能力。}
  \item \textbf{打钩}：在附录 D 的进度表上打钩。这一关过了，整册的乘法就稳了。
\end{enumerate}

\begin{passbox}
\begin{itemize}
  \item 跳着问 20 句，对 18 句以上。
  \item 60 秒对 30 张卡。
  \item 给积能立刻说出两个乘数（如 42 $\to$ 6 和 7）。
\end{itemize}
\end{passbox}
\end{lesson}

\subsection{翻车现场（对着看）}

\begin{center}
\small
\begin{tabularx}{\linewidth}{@{}p{3.8cm}p{4.0cm}X@{}}
\toprule
\textbf{小豆角写出来的样子} & \textbf{其实是想错了什么} & \textbf{豆角爸爸回一句什么} \\
\midrule
七八五十六和六七四十二混 & 数字挨着，靠模糊记忆 & 做成两张卡并排看：42 和 56 差 14。\\
\addlinespace[2pt]
9 的手指法左右错 & 从右手数起 & 一律从左往右，弯指左边当十。\\
\addlinespace[2pt]
只会顺着背，倒着问就卡 & 靠顺序记的 & 抽查必须跳着问、倒着问。\\
\addlinespace[2pt]
忘一句就放弃 & 不会拆数 & 拆成两句会的：$6\times7=30+12$。\\
\addlinespace[2pt]
$9\times6$ 答成 56 & 和 $7\times8$ 混了 & 用 $60-6=54$ 验一遍。\\
\addlinespace[2pt]
表填完就不看了 & 只当练习用 & 问“对角线是什么”“关于什么对称”。\\
\bottomrule
\end{tabularx}
\end{center}

\subsection{抽查：随机开枪（10 题）}

\begin{enumerate}[label=\arabic*.,itemsep=4pt]
  \item 六七 \blank[1.6cm] \qquad 2. 七八 \blank[1.6cm] \qquad 3. 九六 \blank[1.6cm]
  \item $7\times8=$ \blank[1.4cm] \qquad 5. $9\times6=$ \blank[1.4cm] \qquad
  6. $8\times8=$ \blank[1.4cm]
  \item $9\times7=$ \blank[1.4cm] \qquad 8. $6\times9=$ \blank[1.4cm]
  \item $42\div6=$ \blank[1.4cm] \qquad 10. $56\div7=$ \blank[1.4cm]
\end{enumerate}

\begin{ansbox}
1. 四十二\quad 2. 五十六\quad 3. 五十四\quad 4. 56\quad 5. 54\quad
6. 64\quad 7. 63\quad 8. 54\quad 9. 7\quad 10. 8 \\
\textbf{第 9、10 题错}：口诀还不能反着用，先专门练“给积说乘数”。
\end{ansbox}

\subsection{课后野路子}

\begin{itemize}
  \item \textbf{每天 3 分钟抽卡}：抽 20 张，跳出来的立刻翻出来单练。
  \item \textbf{饭前对口令}：豆角爸爸一句小豆角一句，跳着来，两分钟。
  \item \textbf{生活里的 6--9}：一盒鸡蛋 3 排 2 列、一箱牛奶 4 排 3 列、
        一层楼 9 级台阶走 7 层。
  \item \textbf{扑克牌}：抽两张相乘，比谁算得快；三张时先乘再加（3、4、5 $\to 3\times4+5$）。
  \item \textbf{买东西}：一包 8 块，买 7 包多少钱？先口令诀再说钱数。
\end{itemize}
